\subsection{覆叠空间中庞加莱群胚的作用}

设 B 是一个拓扑空间，\(p \colon \mathrm{E} \to \mathrm{B}\) 是满足道路提升性质的分离平展映射（III, p.~302）；如果映射 \(p\) 使 E 成为 B 的覆叠，情形就是如此（上引文）。设 \(b\) 和 \(b'\) 是 B 中的点，\(c \in \Lambda_{b,b'}(\mathrm{B})\) 是 B 中连接 \(b\) 与 \(b'\) 的道路。对纤维 \(x\)、\(\mathrm{E}_b\) 中的每个点 \(x \cdot c\)，记 E 中以 x 为起点且提升 c 的道路的终点为 \(x\)。点 \(c\) 属于纤维 \(x \cdot c\)，并且只依赖于道路 c 的类 \(E_{b'}\)（III, p.~302, 命题 4）；因此将 \(\gamma \in \varpi_{b,b'}(\mathrm{B})\)、\(x \cdot \gamma\) 记为 \(x \cdot c\)。

若 c 是 b 处的常值道路，则有 \(c\)、\(b\)、\(x \cdot \gamma = x\)、\(x\)、\(c\)，因为以 x 为起点的常值道路提升 c。

设 \(b''\) 是 B 中的另一个点，\(c' \in \Lambda_{b',b''}(B)\)。对 \(x\) 中的每个点 \(E_b\)，有：

\[
(x \cdot c) \cdot c ^ {\prime} = x \cdot (c * c ^ {\prime}).\tag{1}
\]

事实上，记 \(\widetilde{c}\) 为以 x 为起点的 c 的提升，记 \(x\)、\(c\)、\(\widetilde{c}'\)、\(c'\) 为以 \(x \cdot c = \widetilde{c}(1)\) 为起点的 c' 的提升。道路 \(\widetilde{c}\) 与 \(\widetilde{c}'\) 可拼接，且 \(\widetilde{c} * \widetilde{c}'\) 是以 x 为起点提升 \(x\)、\(c * c'\) 的道路。

记 \(\varphi_{b,b'}\colon \varpi_{b,b'}(\mathrm{B})\to \mathcal{F}(\mathrm{E}_b;\mathrm{E}_{b'})\) 为满足下列条件的映射：对 \(\gamma \in \varpi_{b,b'}(\mathrm{B})\) 和 \(x\in \mathrm{E}_b\)，有

\[
\varphi_ {b, b ^ {\prime}} (\gamma) (x) = x \cdot \gamma .
\]

对每个 \(\gamma \in \varpi_{b,b'}(B)\) 和每个 \(\gamma' \in \varpi_{b',b''}(B)\)，由关系 (1) 可得

\[
\varphi_ {b, b ^ {\prime \prime}} (\gamma \gamma^ {\prime}) = \varphi_ {b ^ {\prime}, b ^ {\prime \prime}} (\gamma^ {\prime}) \circ \varphi_ {b, b ^ {\prime}} (\gamma).\tag{2}
\]

族 \(\varphi = (\varphi_{b,b'})_{(b,b')\in\mathrm{B}\times\mathrm{B}}\) 因而关于映射 \(\varpi(\mathrm{B})\) 在集合 E 上定义了群胚 \(p\colon\mathrm{E}\to\mathrm{B}\) 的右作用法则（II, p.~167）。称其为与映射 \(\varpi(\mathrm{B})\) 相伴的群胚 \(p\colon\mathrm{E}\to\mathrm{B}\) 的典范作用。映射 \(\varphi_{b,b}\colon\pi_{1}(\mathrm{B},b)\to\mathcal{F}(\mathrm{E}_{b};\mathrm{E}_{b})\) 定义群 \(\pi_{1}(\mathrm{B},b)\) 在纤维 \(\mathrm{E}_{b}\) 上的右作用。称此作用为 \(\pi_{1}(\mathrm{B},b)\) 在纤维 \(\mathrm{E}_{b}\) 上的典范作用。

注. --- 设 \(b\) 和 \(b'\) 是 B 中的两个点，\(x\) 是纤维 \(\mathrm{E}_b\) 中的一点，\(c \in \Lambda_{b,b'}(\mathrm{B})\) 是 B 中的一条道路，\(\widetilde{c}\)、\(x\)、\(c\) 是 E 中以 x 为起点且提升 c 的道路。对所有 \(s \in \mathbf{I}\)，记 \(c_s\) 为道路 \(t \mapsto c(st)\)。E 中以 x 为起点提升 \(x\)、\(c_s\) 的道路是 \(t \mapsto \widetilde{c}(st)\)，其终点是点 \(\widetilde{c}(s)\)。因此，对所有 \(\widetilde{c}(s) = x \cdot c_s\)，有 \(s \in \mathbf{I}\)。

设 B' 是一个拓扑空间，\(p'\colon E'\to B'\) 是满足道路提升性质的平展分离映射。设 \(f\colon B\to B'\) 和 \(g\colon E\to E'\) 是连续映射，满足 \(p'\circ g=f\circ p\)。对 \(b, b'\in B, \gamma\in\varpi_{b,b'}(B)\) 和 \(x\in E_b\)，有：

\[
g (x \cdot \gamma) = g (x) \cdot f _ {*} (\gamma).\tag{3}
\]

事实上，令 c 为 B 中严格同伦类为 \(\gamma\) 的道路；记 \(\widetilde{c}\) 为 E 中以 x 为起点且提升 c 的道路。于是道路 \(g \circ \widetilde{c}\) 是 E' 中以 \(g(x)\) 为起点提升 \(f \circ c\) 的道路。

特别地，当 \(B = B'\) 且 \(f = \operatorname{Id}_{B}\) 时，有

\[
g (x \cdot \gamma) = g (x) \cdot \gamma .\tag{4}
\]

设 \(p \colon \mathrm{E} \to \mathrm{B}\) 和 \(p' \colon \mathrm{E}' \to \mathrm{B}\) 是满足道路提升性质的平展分离映射。设 b 是 B 中的一点。若 \(b\)、\(f \colon \mathrm{E} \to \mathrm{E}'\) 是 B-态射，则映射 \(f_b \colon \mathrm{E}_b \to \mathrm{E}_b'\) 是一个 \(\pi_1(\mathrm{B}, b)\)-态射。

定理 1. --- 设 B 是一个拓扑空间， 是 B 的一个覆叠，b 是 B 中的一点，x 是纤维 \(E_{b}\) 中的一点。

\begin{enumerate}
\def\labelenumi{\alph{enumi})}
\item
  x 在群 \(\pi_{1}(\mathrm{B},b)\) 的典范作用下的轨道等于 \(E_{b}\) 与 x 在 E 中的道路连通分支的交。特别地，若空间 E 是道路连通的，则此作用是传递的。
\item
  x 的固定子（A, I, p.~51）是 \(p_{*}(\pi_{1}(E,x))\) 的子群 \(\pi_{1}(B,b)\)。
\item
  对每个 \(\gamma \in \pi_1(\mathrm{B},b)\)，有 \(p_*(\pi_1(\mathrm{E},x)) = \operatorname{Int}(\gamma)(p_*(\pi_1(\mathrm{E},x \cdot \gamma)))\)。
\item
  如果空间 B 是连通的且覆叠 E 是伽罗瓦覆叠，则子群 \(p_{*}(\pi_{1}(\mathrm{E},x))\) 在 \(\pi_{1}(\mathrm{B},b)\) 中是正规子群。
\end{enumerate}

断言 a) 是直接的。断言 b) 由命题 4 的推论 2（III, p.~303）得到。断言 c) 由 b) 以及 A, I, p.~52 的命题 2 得到。最后，假设 B 是连通的，且 E 是 B 的伽罗瓦覆叠。对每个 \(\gamma \in \pi_1(B,b)\)，存在 E 的一个 B-自同构 \(g\)，使得 \(g(x) = x \cdot \gamma\)（I, p.~102, 定理 2, c)）。有 \(p = p \circ g\)，从而 \(p_*(\pi_1(E,x)) = p_*(\pi_1(E,x \cdot \gamma))\)。因此断言 d) 由 c) 得到。
% semantic-fragments: legacy-input-seam
\relax{}
